% ============================================================
% tree.py paper -- Section 1 + Section 2
% Wiley NJD two-column format, 80-char line limit
% ============================================================
% 

% We mar18
% early: vavetas: not all ai can be repal;ced by these techniquens. dialgo generation stays, itnefacing begietween disptierat tools, feature extarction over ehtregenous sources, covnersion to differenr report forsts,  all remains, con

% but one to look at this is what arethe tools that udner the hood an agentic system will ahave to use. soemthings these will be compelx embeddigns and sometimes they wont. and when they arent' systemrun faste with lset cost to the suer


% sstartd obejctslions:
% - simple is ok but ore siphisticated sytems work better. to wichw e reply, ahve you tied? and useds stats to comaprereuslts before afetr? you sure that the new says are stsitically distinguishable by more than a smalle ffect? we ask since as shown here, we ae apply such stats, we rarely see an improvement (which of course ight not be true for other domaisn)

% - only for simple data. jsut low diemansiaon talbles tables (1000 colsumsn or less). well inreply we say "jsut" tables is not really just. for jighly structure soruces, LLMs are known to ahv troube. also, we have results from high diesnioanl text minign data wehre we kick ass

% section: run faster

% 1. runtimes>
% 2. labelling effort

%
\documentclass[AMA,STIX2COL]{WileyNJDv5}

\begin{document}

\title{The \texttt{EZR.py} Toolkit: How Far Can We Go
       With Simple AI?}
\author{Anonymous Authors}

\maketitle

% ============================================================
\section{Introduction}\label{sec:intro}
% ============================================================

Software engineering is, for the most part, racing toward
complexity. Few-shot prompting has evolved into agentic
workflows where large language models~(LLMs) orchestrate
multi-agent collaboration and repository-scale
comprehension. There is much success with this paradigm,
and much promise for what such systems will
deliver~\cite{ma2023scope,wang2022machine,%
Dong2025SurveyCodeAgents,DBLP:journals/corr/abs-2507-15003}.


Yet there are growing reasons to reconsider that ``complexity-first'' approach.

Sustainability
concerns loom large: if LLM services are no longer
venture-subsidized, will their costs rise by one to two
orders of magnitude? And quite apart from economics, a
long string of open challenges persists (Table~\ref{cautions}).
Problems with explanation make certification and ongoing
validation difficult, particularly for high-stakes
applications in law, medicine, and safety-critical systems.
Not all deployment environments are CPU-intensive, either;
edge computing demands something leaner.

So while we say ``agentic, yes,'' we ask whether every
\emph{part} of an agentic pipeline must be a
resource-intensive LLM. Recent results offer a reason for
optimism. PromiseTune~\cite{chen2026promisetune} shows that
configuration landscapes collapse into causally promising
regions; the BINGO effect~\cite{ganguly2025lowgodatalightse}
shows that SE data concentrates in surprisingly few
multi-dimensional buckets. Both results suggest that the
effective state space of many SE problems is far smaller
than it appears---and that a handful of lightweight probes
can navigate it.

To test that hypothesis we built \texttt{tree.py}: a
symbolic AI written with the simplest methods we could
muster, then extended with additional learners to measure
how little new code each new capability required. The usual
concern is that simpler methods perform worse. To stress-test
that concern we evaluated on 120+ SE tasks from the MOOT
repository~\cite{menzies2026moot}---up to 1,000 independent
variables, up to 8 goals, and some datasets with hundreds
of thousands of examples. The verdict: our simpler methods
are remarkably competitive. They require very few labeled
examples, complete in seconds, produce insightful
explanations of complex problems, and fit in a $\sim$300-line
core plus $\sim$60-line application modules.

Caveats are necessary. We do not claim that every SE task
can be simplified. Some inherently require LLMs: translation,
text summarization, dialog generation, extracting structure
from large heterogeneous sources. What we show is that
\emph{some} SE tasks---and the list is not short---can be
solved without trillion-parameter models. The MOOT benchmark
spans regression, prediction, multi-objective optimization,
configuration, performance tuning, hyperparameter
optimization, product line engineering, project health
forecasting, defect prediction, software testing, process
and cost estimation, cross-domain generalization, data
generation, and text mining. Clearly there is a large space
of possible simplifications, and future work will be needed
to see how far simplicity can reach.

\begin{table}[!t]
\caption{Some reported issues with LLMs.}\label{cautions}
\begin{tabular}{l|p{5cm}}
\textbf{Issue} & \textbf{Evidence} \\ \hline
Hallucinations
  & Over half of ChatGPT answers contain incorrect
    information~\cite{S12} \\
Human error
  & Humans overlook AI errors 39\% of the time~\cite{S12} \\
Code quality
  & 70\% faulty: 29\% correct, 51\% partial,
    20\% wrong~\cite{S15} \\
Expert caution
  & Google developers accept $\le 25\%$ of AI
    suggestions~\cite{tabachnyk22completion} \\
Training bias
  & LLMs regenerate bugs from training
    data~\cite{S13} \\
Energy         & High consumption. \\
Explanation
  & Subsymbolic: no causal or explanatory
    reasoning; hard to debug or edit. \\
Verification   & Prompt output is rarely verified. \\
Planning
  & ``Vibe coding'' skips planning, raising
    test costs. \\
Ownership
  & Copy-pasting reduces developer
    job satisfaction. \\
Bubble bursting & Withdrawal of investor capital? \\
Subsidies ending?
  & If capital moves, will LLM costs
    rise 10$\times$--100$\times$? \\
Cognitive debt
  & Agentic AI lets teams move fast while quietly
    eroding shared understanding of system
    design~\cite{CognitiveDebt} \\
Maintenance
  & Lack of ``under-the-hood'' skills
    $\rightarrow$ disasters? \\
\end{tabular}
\end{table}

% ============================================================
\section{Why Simplicity?}\label{sec:simplicity}
% ============================================================

\subsection{The Case for Simple Methods}

Simplicity is no new idea. Ptolemy~(130~A.D.) urged
explaining phenomena by ``the simplest hypothesis
possible.'' William of Occam warned that entities must
not be multiplied beyond necessity. The documented benefits
are concrete: lower compute and energy
costs~\cite{calrero15,monserrate2022cloud}, easier
explanation~\cite{veerappa2011understanding}, and faster
tuning~\cite{fu2016tuning}. Rudin~\cite{rudin2019stop}
puts it plainly:

\begin{quote}
{\em ``Stop explaining black box machine learning models
for high stakes decisions and use interpretable models
instead.''}
\end{quote}

History bears this out. In the ``Shockingly Simple'' paper,
Menzies~\cite{Menzies21shockingly} documents a recurring
pattern: a few variables (``keys'') govern the rest of a
system, so that very simple models---built from small
samples, greedy search, or Na{\"i}ve Bayes---routinely
match or beat state-of-the-art performance while running
100$\times$--1000$\times$
faster~\cite{majumder2018500+,grinsztajn2022why,%
Tawosi23,Ling24,Fu17}. Table~\ref{less} catalogues
illustrative examples from the literature.

\begin{table}[!t]
\caption{Selected cases where simpler methods matched
  or outperformed complex ones.}\label{less}
\begin{tabular}{l|p{5cm}}
\textbf{Domain} & \textbf{Finding} \\ \hline
Effort estimation
  & SVM+TFIDF beat deep learners and ran
    100$\times$ faster~\cite{Tawosi23} \\
Defect prediction
  & K-Means pre-clustering matched deep learning
    at 500$\times$ speedup~\cite{majumder2018500+} \\
Data synthesis
  & Parametric statistics equalled GANs with
    less compute~\cite{Ling24} \\
Hyperparameter opt.
  & Very fast tuning worked as well as
    complex grid search~\cite{Fu17} \\
Tabular data
  & Tree-based models outperform deep nets
    on medium-sized tables~\cite{grinsztajn2022why} \\
\end{tabular}
\end{table}

Yet the lesson goes unheeded. A survey of 229 SE papers
using LLMs found that only 5\% compared against a simpler
baseline~\cite{Hou24}. We conjecture a cognitive cause:
Adams and Fillon et al.~\cite{adams2021people,Fillon25}
found that across 400+ subjects given design tasks, people
chose additive over subtractive changes at a 4:1 ratio.
{\em Humans needlessly complicate their designs.}

\subsection{The BINGO Effect: Encouragement to Simplify}

The ``How Low Can You Go?'' data-light SE
challenge~\cite{ganguly2025lowgodatalightse} provides
empirical grounding for our optimism. Across 127 SE
optimization tasks from the MOOT repository, the
\emph{BINGO effect} holds: when data is partitioned into
$n$ buckets across $d$ dimensions divided into $b$ bins,
most data concentrates in surprisingly few buckets
($n \ll b^d$). In one representative study with 10,000
rows, the expected $b^d = 4{,}096$ buckets collapsed to
just 100 occupied ones. Because the effective search space
is so small, a handful of labeled examples is enough to
find near-optimal solutions.

Concretely, using only a few dozen labels, simple
pool-based methods---diversity sampling, a minimal Bayesian
learner, or even random probes---achieved over 90\% of the
best results found by state-of-the-art optimizers like
SMAC, TPE, and DEHB, while running more than two orders
of magnitude faster. This is a profound encouragement to
see how far we can go with very simple AI.

\subsection{Our Test Bed: The MOOT Repository}

To put these ideas to work we used the MOOT repository of
multi-objective optimization
tasks~\cite{menzies2026moot}. MOOT's 120+ tasks span a
wide range of SE challenges; Table~\ref{moot} summarises
the diversity. Each task is a table of $x$ inputs and $y$
goals; goal columns are annotated with $+/-$ to indicate
maximize/minimize. Dataset sizes range from 93 to
166,975 rows; dimensionality ranges from 3 to 1,044
independent variables; goal counts range from 1 to 8.
This variety makes MOOT a demanding proving ground for
any claim that simple methods suffice.

One striking finding from our experience with MOOT is
that the inference machinery is largely shared across
tasks. For example:
\begin{itemize}
\item Coding a nearest-neighbor classifier provided nearly
  everything needed for stochastic optimization.
\item Coding an optimizer provided nearly everything needed
  for software configuration and tuning.
\end{itemize}
This convergence suggests a radical reframing: rather than
a zoo of specialized algorithms, AI may reduce to a
handful of core operations, each ``different'' technique
merely a thin variation on the same kernel. The field has
been making things harder than they need to be.
\texttt{tree.py} is an existence proof that a simpler
foundation is both possible and practical.

\bibliography{main}
\end{document}